% GENERATED FILE. Edit canonical lessons, then run npm run generate:books.
% canonical-source-hash: fnv1a64:165f26903ada2abc
% canonical-lessons: PA-C302-kiari, PA-C302-ruri, PA-C302-khilarna, PA-C302-uggna, PA-C302-sarhom

\chapter{Sowing Time}
\label{ch:pa-a2-sowing-time}
\hlchaptermodality{302}

\begin{chapteropening}
\textbf{By the end of this chapter:} \emph{I can talk about sowing and sprouting.}

Complete the last lesson of chapter 302: I can talk about sowing and sprouting.
\end{chapteropening}

\section[\texorpdfstring{\pa{ਕਿਆਰੀ} (kiārī)}{ਕਿਆਰੀ (kiārī)}]{\pa{ਕਿਆਰੀ} (kiārī) — a garden bed; a bed in a field}

\label{lesson:PA-C302-kiari}



\begin{quote}
Before the new one: say the Punjabi for to plough, then the Punjabi for to yoke.
\end{quote}

\subsection*{You'll want to know: \pa{ਕਿਆਰੀ}}
\textbf{\pa{ਕਿਆਰੀ}} — \emph{kiārī} — \textquotedblleft{}a garden bed; a bed in a field\textquotedblright{}.

\begin{quote}
\pa{ਮਾਤਾ} \pa{ਜੀ} \pa{ਕਿਆਰੀ} \pa{ਵਿੱਚ} \pa{ਸਾਗ} \pa{ਬੀਜਦੇ} \pa{ਹਨ।}
\end{quote}

\begin{quote}
\emph{mātā jī kiārī vich sāg bījde han.}
\end{quote}

\begin{quote}
Mother sows greens in the garden bed.
\end{quote}

A feminine noun.

\begin{cousinweb}
\textbf{\pa{ਕਿਆਰੀ}} (\emph{kiārī}) is inherited from Sanskrit \emph{kēdāra}, \textquotedblleft{}field\textquotedblright{}; CDIAL lists Punjabi \emph{kiārī}, \textquotedblleft{}bed in a field or garden\textquotedblright{}, under it (CDIAL 3463).
\end{cousinweb}

\begin{grammarlens}[title={What you've built}]
One more word for this chapter.
\end{grammarlens}

\subsection*{Your turn}
\begin{itemize}
\raggedright
  \item \emph{Say it:} \emph{kiārī}
  \item \emph{Say it:} \emph{kiārī}, once more
  \item \emph{Say it:} \emph{mātā jī kiārī vich sāg bījde han.} — Mother sows greens in the garden bed.
  \item \emph{Recall:} say the Punjabi for to plough, then the Punjabi for to yoke, then say \emph{kiārī} again
\end{itemize}

\begin{tcolorbox}[breakable,colback=teal!4,colframe=teal!35!black,title={Before you move on}]
What does \pa{ਕਿਆਰੀ} mean? (\textquotedblleft{}A garden bed; a bed in a field\textquotedblright{}.) Say it once more.
\end{tcolorbox}
\section[\texorpdfstring{\pa{ਰੂੜੀ} (rūṛī)}{ਰੂੜੀ (rūṛī)}]{\pa{ਰੂੜੀ} (rūṛī) — a manure heap}

\label{lesson:PA-C302-ruri}



\begin{quote}
Before the new one: say the Punjabi for to yoke, then the Punjabi for a garden bed.
\end{quote}

\subsection*{You'll want to know: \pa{ਰੂੜੀ}}
\textbf{\pa{ਰੂੜੀ}} — \emph{rūṛī} — \textquotedblleft{}a manure heap\textquotedblright{}.

\begin{quote}
\pa{ਰੂੜੀ} \pa{ਪਿੰਡ} \pa{ਦੇ} \pa{ਬਾਹਰ} \pa{ਹੈ।}
\end{quote}

\begin{quote}
\emph{rūṛī piṇḍ de bāhar hai.}
\end{quote}

\begin{quote}
The manure heap is outside the village.
\end{quote}

A feminine noun.

\begin{cousinweb}
\textbf{\pa{ਰੂੜੀ}} (\emph{rūṛī}) is inherited from Sanskrit \emph{rūḍhi}, \textquotedblleft{}rise, increase\textquotedblright{}; CDIAL lists Punjabi \emph{rūṛī}, \textquotedblleft{}heap of manure outside village\textquotedblright{}, under it (CDIAL 10802).
\end{cousinweb}

\begin{grammarlens}[title={What you've built}]
One more word for this chapter.
\end{grammarlens}

\subsection*{Your turn}
\begin{itemize}
\raggedright
  \item \emph{Say it:} \emph{rūṛī}
  \item \emph{Say it:} \emph{rūṛī}, once more
  \item \emph{Say it:} \emph{rūṛī piṇḍ de bāhar hai.} — The manure heap is outside the village.
  \item \emph{Recall:} say the Punjabi for to yoke, then the Punjabi for a garden bed, then say \emph{rūṛī} again
  \item \emph{Write it:} copy \textbf{\pa{ਰੂੜੀ}} once, checking the whole word
\end{itemize}

\begin{tcolorbox}[breakable,colback=teal!4,colframe=teal!35!black,title={Before you move on}]
What does \pa{ਰੂੜੀ} mean? (\textquotedblleft{}A manure heap\textquotedblright{}.) Say it once more.
\end{tcolorbox}
\section[\texorpdfstring{\pa{ਖਿਲਾਰਨਾ} (khilārnā)}{ਖਿਲਾਰਨਾ (khilārnā)}]{\pa{ਖਿਲਾਰਨਾ} (khilārnā) — to scatter; to spread about}

\label{lesson:PA-C302-khilarna}



\begin{quote}
Before the new one: say the Punjabi for a garden bed, then the Punjabi for a manure heap.
\end{quote}

\subsection*{You'll want to know: \pa{ਖਿਲਾਰਨਾ}}
\textbf{\pa{ਖਿਲਾਰਨਾ}} — \emph{khilārnā} — \textquotedblleft{}to scatter; to spread about\textquotedblright{}.

\begin{quote}
\pa{ਮੈਂ} \pa{ਖੇਤ} \pa{ਵਿੱਚ} \pa{ਬੀਜ} \pa{ਖਿਲਾਰਦਾ} \pa{ਹਾਂ।}
\end{quote}

\begin{quote}
\emph{maiṁ khet vich bīj khilārdā hāṁ.}
\end{quote}

\begin{quote}
I scatter seed in the field.
\end{quote}

Past (I …): \pa{ਮੈਂ} \pa{ਬੀਜ} \pa{ਖਿਲਾਰੇ}, \emph{maiṁ bīj khilāre}.

\begin{cousinweb}
\textbf{\pa{ਖਿਲਾਰਨਾ}} (\emph{khilārnā}) is inherited from a reconstructed \emph{*khiḍ}, \textquotedblleft{}be open, expand\textquotedblright{}; CDIAL lists Punjabi \emph{khillarṇā}, \textquotedblleft{}to be scattered\textquotedblright{}, and its causative \emph{khilāraṇā} under it (CDIAL 3882).
\end{cousinweb}

\begin{grammarlens}[title={What you've built}]
One more word for this chapter.
\end{grammarlens}

\subsection*{Your turn}
\begin{itemize}
\raggedright
  \item \emph{Say it:} \emph{khilārnā}
  \item \emph{Say it:} \emph{khilārnā}, once more
  \item \emph{Say it:} \emph{maiṁ khet vich bīj khilārdā hāṁ.} — I scatter seed in the field.
  \item \emph{Recall:} say the Punjabi for a garden bed, then the Punjabi for a manure heap, then say \emph{khilārnā} again
  \item \emph{Write it:} copy \textbf{\pa{ਖਿਲਾਰਨਾ}} once, checking the whole word
\end{itemize}

\begin{tcolorbox}[breakable,colback=teal!4,colframe=teal!35!black,title={Before you move on}]
What does \pa{ਖਿਲਾਰਨਾ} mean? (\textquotedblleft{}To scatter; to spread about\textquotedblright{}.) Say it once more.
\end{tcolorbox}
\section[\texorpdfstring{\pa{ਉੱਗਣਾ} (uggṇā)}{ਉੱਗਣਾ (uggṇā)}]{\pa{ਉੱਗਣਾ} (uggṇā) — to sprout; to come up}

\label{lesson:PA-C302-uggna}



\begin{quote}
Before the new one: say the Punjabi for a manure heap, then the Punjabi for to scatter.
\end{quote}

\subsection*{You'll want to know: \pa{ਉੱਗਣਾ}}
\textbf{\pa{ਉੱਗਣਾ}} — \emph{uggṇā} — \textquotedblleft{}to sprout; to come up\textquotedblright{}.

\begin{quote}
\pa{ਮੀਂਹ} \pa{ਤੋਂ} \pa{ਬਾਅਦ} \pa{ਬੀਜ} \pa{ਉੱਗਣਗੇ।}
\end{quote}

\begin{quote}
\emph{mīṁh ton bāad bīj uggaṇge.}
\end{quote}

\begin{quote}
After the rain the seeds will sprout.
\end{quote}

Past (it …): \pa{ਉਹ} \pa{ਘਾਹ} \pa{ਉੱਗਿਆ}, \emph{oh ghāh uggiā}.

\begin{cousinweb}
\textbf{\pa{ਉੱਗਣਾ}} (\emph{uggṇā}) is inherited from a reconstructed \emph{*udgāti}, \textquotedblleft{}rises, comes forth\textquotedblright{}; CDIAL lists Punjabi \emph{uggṇā}, \textquotedblleft{}to spring up, grow\textquotedblright{}, under it (CDIAL 1954).
\end{cousinweb}

\begin{grammarlens}[title={What you've built}]
One more word for this chapter.
\end{grammarlens}

\subsection*{Your turn}
\begin{itemize}
\raggedright
  \item \emph{Say it:} \emph{uggṇā}
  \item \emph{Say it:} \emph{uggṇā}, once more
  \item \emph{Say it:} \emph{mīṁh ton bāad bīj uggaṇge.} — After the rain the seeds will sprout.
  \item \emph{Recall:} say the Punjabi for a manure heap, then the Punjabi for to scatter, then say \emph{uggṇā} again
  \item \emph{Write it:} copy \textbf{\pa{ਉੱਗਣਾ}} once, checking the whole word
\end{itemize}

\begin{tcolorbox}[breakable,colback=teal!4,colframe=teal!35!black,title={Before you move on}]
What does \pa{ਉੱਗਣਾ} mean? (\textquotedblleft{}To sprout; to come up\textquotedblright{}.) Say it once more.
\end{tcolorbox}
\section[\texorpdfstring{\pa{ਸਰ੍ਹੋਂ} (sarhoṁ)}{ਸਰ੍ਹੋਂ (sarhoṁ)}]{\pa{ਸਰ੍ਹੋਂ} (sarhoṁ) — mustard (the crop)}

\label{lesson:PA-C302-sarhom}



\begin{quote}
Before the new one: say the Punjabi for to scatter, then the Punjabi for to sprout.
\end{quote}

\subsection*{You'll want to know: \pa{ਸਰ੍ਹੋਂ}}
\textbf{\pa{ਸਰ੍ਹੋਂ}} — \emph{sarhoṁ} — \textquotedblleft{}mustard (the crop)\textquotedblright{}.

\begin{quote}
\pa{ਮਾਘ} \pa{ਵਿੱਚ} \pa{ਸਰ੍ਹੋਂ} \pa{ਦੇ} \pa{ਫੁੱਲ} \pa{ਪੀਲੇ} \pa{ਹੁੰਦੇ} \pa{ਹਨ।}
\end{quote}

\begin{quote}
\emph{māgh vich sarhoṁ de phull pīle hunde han.}
\end{quote}

\begin{quote}
In Magh the mustard flowers are yellow.
\end{quote}

A feminine noun.

\begin{cousinweb}
\textbf{\pa{ਸਰ੍ਹੋਂ}} (\emph{sarhoṁ}) is inherited from Sanskrit \emph{sarṣapa}, \textquotedblleft{}mustard seed\textquotedblright{}; CDIAL lists Punjabi \emph{sarhõ}, \textquotedblleft{}rape seed\textquotedblright{}, under it (CDIAL 13281).
\end{cousinweb}

\begin{grammarlens}[title={What you've built}]
One more word for this chapter.
\end{grammarlens}

\subsection*{Your turn}
\begin{itemize}
\raggedright
  \item \emph{Say it:} \emph{sarhoṁ}
  \item \emph{Say it:} \emph{sarhoṁ}, once more
  \item \emph{Say it:} \emph{māgh vich sarhoṁ de phull pīle hunde han.} — In Magh the mustard flowers are yellow.
  \item \emph{Recall:} say the Punjabi for to scatter, then the Punjabi for to sprout, then say \emph{sarhoṁ} again
  \item \emph{Write it:} copy \textbf{\pa{ਸਰ੍ਹੋਂ}} once, checking the whole word
\end{itemize}

\begin{tcolorbox}[breakable,colback=teal!4,colframe=teal!35!black,title={Before you move on}]
What does \pa{ਸਰ੍ਹੋਂ} mean? (\textquotedblleft{}Mustard (the crop)\textquotedblright{}.) Say it once more.
\end{tcolorbox}
